\documentclass[11pt,a4paper]{article}
\usepackage{brainteaser}

\begin{document}
\brainteasertitle{18}{Les dés de Sicherman}

\begin{statementbox}
On lance deux dés ordinaires équilibrés à six faces, numérotées de $1$ à $6$.

On souhaite les remplacer par deux autres dés équilibrés à six faces. Les
faces peuvent porter des nombres entiers strictement positifs répétés, mais
chaque dé doit toujours posséder exactement six faces.

Peut-on numéroter ces deux dés autrement qu'avec
$1,2,3,4,5,6$, tout en conservant exactement la même distribution des sommes?
Si oui, déterminer les nombres inscrits sur les faces des deux dés.
\end{statementbox}

\begin{hintsbox}
\begin{enumerate}
  \item Associer à un dé dont les faces sont $d_1,\ldots,d_6$ le polynôme
  $x^{d_1}+\cdots+x^{d_6}$.
  \item Que représente le coefficient de $x^s$ dans le produit des
  polynômes associés aux deux dés?
  \item Le polynôme d'un dé ordinaire est
  \[
  D(x)=x+x^2+\cdots+x^6.
  \]
  Factoriser $D(x)$ dans $\mathbb Z[x]$.
  \item Si $A(x)$ et $B(x)$ représentent les dés recherchés, utiliser
  $A(x)B(x)=D(x)^2$ et $A(1)=B(1)=6$.
  \item Parmi les facteurs de $D(x)^2$, lequel peut être réparti différemment
  entre $A(x)$ et $B(x)$ sans modifier leur valeur en $x=1$?
\end{enumerate}
\end{hintsbox}

\correction

Pour un dé, notons $a_k$ le nombre de faces portant l'entier $k$ et
introduisons le polynôme générateur
\[
A(x)=\sum_{k\geq1}a_kx^k.
\]
Le coefficient de $x^k$ est donc le nombre de faces portant $k$, et
$A(1)=6$ puisque le dé possède six faces.

Pour deux dés représentés par $A(x)$ et $B(x)$, le coefficient de $x^s$ dans
$A(x)B(x)$ vaut
\[
\sum_{i+j=s}a_i b_j.
\]
Il compte exactement les couples de faces dont la somme est $s$. Les deux dés
recherchés ont donc la même distribution des sommes que deux dés ordinaires
si et seulement si
\[
A(x)B(x)=D(x)^2,
\qquad
D(x)=x+x^2+x^3+x^4+x^5+x^6.
\]

\subsection*{Factorisation du dé ordinaire}

On factorise
\[
\begin{aligned}
D(x)
&=x(1+x+x^2+x^3+x^4+x^5)\\
&=x(x+1)(x^2+x+1)(x^2-x+1).
\end{aligned}
\]
Ainsi,
\[
D(x)^2
=x^2(x+1)^2(x^2+x+1)^2(x^2-x+1)^2.
\]

Les nombres inscrits étant strictement positifs, $A$ et $B$ sont tous deux
divisibles par $x$. Comme le coefficient de $x^2$ dans leur produit vaut $1$,
chacun contient exactement un facteur $x$ et possède exactement une face
portant le nombre $1$.

Les trois autres facteurs affichés sont irréductibles dans $\mathbb Z[x]$.
L'unicité de la factorisation impose donc, pour certains exposants
$\alpha,\beta,\gamma\in\{0,1,2\}$,
\[
A(x)=x(x+1)^\alpha(x^2+x+1)^\beta(x^2-x+1)^\gamma,
\]
les facteurs complémentaires appartenant à $B(x)$.

Comme $A(1)=6$, on a
\[
6=A(1)=2^\alpha3^\beta.
\]
La seule possibilité est
\[
\alpha=\beta=1.
\]
En revanche, $x^2-x+1$ vaut $1$ en $x=1$: son exposant $\gamma$ peut valoir
$0$, $1$ ou $2$ sans changer le nombre de faces.

\subsection*{Les différentes répartitions}

Si $\gamma=1$, alors
\[
A(x)=B(x)=D(x),
\]
et l'on retrouve simplement les deux dés ordinaires.

Prenons maintenant $\gamma=0$. Le premier polynôme devient
\[
\begin{aligned}
A(x)
&=x(x+1)(x^2+x+1)\\
&=x+2x^2+2x^3+x^4.
\end{aligned}
\]
Le premier dé porte donc les nombres
\[
1,\quad2,\quad2,\quad3,\quad3,\quad4.
\]

Le second reçoit les deux facteurs restants $x^2-x+1$:
\[
\begin{aligned}
B(x)
&=x(x+1)(x^2+x+1)(x^2-x+1)^2\\
&=x+x^3+x^4+x^5+x^6+x^8.
\end{aligned}
\]
Ses faces portent donc
\[
1,\quad3,\quad4,\quad5,\quad6,\quad8.
\]

Le cas $\gamma=2$ échange seulement les rôles des deux dés. Il ne fournit pas
une nouvelle paire.

\begin{resultbox}
À l'ordre des dés près, l'unique solution non standard en entiers strictement
positifs est
\[
\boxed{(1,2,2,3,3,4)}
\qquad\text{et}\qquad
\boxed{(1,3,4,5,6,8)}.
\]
Ces deux dés reproduisent exactement la distribution des sommes de deux dés
ordinaires.
\end{resultbox}

\end{document}
